\documentclass[12pt]{article}
\usepackage[utf8]{inputenc}
\usepackage[T1]{fontenc}
\usepackage{hyperref}
\title{Kodeoha Language Resources for K-12 Curriculum Development}
\author{Compiled by Claude Sonnet 5 (claude-sonnet-5) for Andr\'as Kornai}
\date{}
\begin{document}
\maketitle
\section{Introduction}
Kodeoha (also Kondeha, Kondea; ISO 639-3 vko) is an Austronesian language of the Bungku-Tolaki group, spoken near Lasusua, North Kolaka regency, Southeast Sulawesi, Indonesia. Estimates are about 1,500 speakers (1999) and a community of 2,200. Endangered Languages rates it threatened and notes that no literacy materials exist; speakers also use Bugis and Indonesian.
\section{Primary Educational Resources}
None found.
\section{Linguistic and Reference Materials}
\begin{itemize}
\item Wikipedia entry (mentions David Mead's 1999 work on Bungku-Tolaki languages): \url{https://en.wikipedia.org/wiki/Kodeoha_language}
\item Endangered Languages Project entry and sample: \url{https://www.endangeredlanguages.com/lang/1899/samples/11628}
\end{itemize}
\section{Religious and Community Resources}
Joshua Project people-group profile: \url{https://joshuaproject.net/people_groups/12733}
\section{Digital and Audio Resources}
Global Recordings Network language page: \url{https://globalrecordings.net/en/language/vko}
\section{Curriculum Development Context}
No existing curriculum, orthography materials or school use were found. Development would start from documentation and community work, with Indonesian as the likely school language.
\end{document}